\chapter{Critical Results --- How Laboratories Escalate Urgent Findings}
\label{app:critical_values}

\clearpage
\section{What Is a Critical Result?}

A critical result is a laboratory result that a laboratory has designated for urgent communication because, in the relevant clinical context, it may indicate a serious and time-sensitive condition. Critical limits are not universal reference ranges. They are established and periodically reviewed by each laboratory or health system with consideration of the assay, specimen type, patient population, clinical setting, and applicable accreditation or regulatory requirements.

The same numerical result can require different action in different settings. A result may also be spurious because of hemolysis, contamination, delayed processing, an incorrect specimen, or an analytical interference. A critical result should therefore be verified and interpreted with the patient's symptoms, history, medications, vital signs, previous results, and related tests.

\begin{notebox}[Important]
This appendix is educational and does not provide a universal critical-value table. Use the reporting laboratory's current critical limits and notification procedure. Patients should follow the instructions from their clinician or laboratory; severe symptoms such as chest pain, severe shortness of breath, confusion, seizure, fainting, or uncontrolled bleeding require emergency care regardless of whether a laboratory result is available.
\end{notebox}

\section{How a Laboratory Handles a Critical Result}

The usual workflow is:
\begin{enumerate}
    \item Check specimen identity, integrity, collection conditions, and analytical flags.
    \item Repeat or confirm the measurement when required by the local procedure, especially when a pre-analytical problem or unexpected change is suspected.
    \item Notify the responsible clinician or authorized care team using the laboratory's defined escalation pathway and document the communication.
    \item The clinical team assesses the patient and determines treatment, repeat testing, monitoring, or referral. A laboratory notification is not itself a diagnosis or a treatment order.
\end{enumerate}

\section{Examples of Results That Commonly Require Urgent Assessment}

The following are clinical examples, not universal thresholds:
\begin{itemize}
    \item Markedly abnormal potassium, sodium, calcium, magnesium, or bicarbonate, particularly with arrhythmia, weakness, seizure, altered mental status, or severe dehydration.
    \item Severe hypoglycemia or marked hyperglycemia with symptoms, ketones, acidosis, dehydration, or altered mental status.
    \item A markedly low hemoglobin or platelet count, severe neutropenia with fever, or a major unexpected change in a blood count.
    \item Markedly abnormal coagulation results in a patient who is bleeding, has a high-risk procedure, or is receiving anticoagulant therapy.
    \item A substantially elevated cardiac troponin with symptoms, ECG changes, or a significant rise or fall. Troponin indicates myocardial injury; it does not independently diagnose myocardial infarction or determine whether catheterization is required.
    \item Severe acidemia, alkalemia, hypoxemia, or hypercapnia, interpreted with the specimen source, oxygen therapy, ventilation, chronicity, and clinical condition.
    \item A potentially toxic drug or toxin concentration. Toxicology results must be interpreted using the substance, formulation, time since exposure, symptoms, organ function, and poison-center or toxicology guidance; isolated numbers are not universal treatment thresholds.
\end{itemize}

\section{Patient Guidance}

If a clinician or laboratory contacts a patient about a critical result, the patient should follow the stated instructions and should not independently change medication doses. If the responsible clinician cannot be reached and the patient has severe or rapidly worsening symptoms, the patient should seek emergency care. A result copied from a portal should be discussed with the ordering clinician because the reporting laboratory's reference interval, critical limit, specimen comments, and clinical context are essential.

\section{Why Fixed Tables Are Unsafe}

Reference intervals and critical limits differ between laboratories and methods. Examples include high-sensitivity troponin assays with different upper reference limits, blood-gas values that depend on arterial or venous sampling and oxygen support, and drug concentrations whose interpretation depends on dose timing and pharmacokinetics. A fixed table in a general reference cannot safely replace the local laboratory's current policy.

\section{Related Information}

For interpretation, consult the individual test chapter, the laboratory report, and current institutional protocols. For suspected poisoning, contact the regional poison center or a medical toxicologist. For suspected sepsis, acute coronary syndrome, severe electrolyte disturbance, or respiratory failure, follow the relevant emergency and critical-care protocol.

\clearpage
